\documentclass[a4paper,10pt,headings=small]{scrartcl}

\usepackage{amsmath,amssymb,amsthm}
\usepackage[american]{babel} % American English
\usepackage[utf8x]{inputenc}
\usepackage{xr}
%\usepackage{xr-hyper} 
\usepackage{hyperref} % hyperlinks
  \hypersetup{hidelinks}

\usepackage[capitalize,noabbrev]{cleveref}
  \creflabelformat{equation}{#2#1#3}
\theoremstyle{definition}
\newtheorem{defn}{Definition}
\newtheorem{exmp}{Example}
\newtheorem{prop}{Property}
\newtheorem{remk}{Remark}

% other packages
\usepackage{algorithmicx}
\usepackage[numbers]{natbib}
\usepackage{csquotes,url}%natbib,
\usepackage{booktabs,graphicx,siunitx,tikz}
\usepackage[figuresright]{rotating}
\usepackage[official]{eurosym}
\usepackage{comment}

% External document reference
\externaldocument{recupred} 

% Title and author
\title{Appendix: The Power of Linear Recurrent Neural Networks}
\author{Frieder Stolzenburg, Sandra Litz, Olivia Michael, and Oliver Obst}
\date{}

\begin{document}

\maketitle

\section{Proof of \cref{rem}}\label{proof:rem}

\noindent For a function $x$ and its first derivative $\dot{x}$ with
respect to time $t$, we have
\[
  \dot{x}(t) = \lim\limits_{\tau \to 0} \frac{x(t+\tau)-x(t)}{\tau}
	\quad\text{and hence}\quad
  x(t+\tau) \approx x(t) + \tau\,\dot{x}(t)
\]
for small time steps $\tau>0$. We can apply this to $x^{(k)}(t)$ for $k \ge 0$
in the difference of \cref{diff} between the times $t+\tau$ and $t$ divided by
$\tau$ and obtain:
\begin{eqnarray*}
  0 & = & \frac{\displaystyle\sum_{k=0}^n c_k\,x^{(k)}(t+\tau)-\sum_{k=0}^n c_k\,x^{(k)}(t)}{\tau}\\
    & = & \sum_{k=0}^{n-1} c_k \frac{x^{(k)}(t+\tau)-x^{(k)}(t)}{\tau} + c_n \frac{x^{(n)}(t+\tau)-x^{(n)}(t)}{\tau}\\
    & \approx & \sum_{k=0}^{n-1} c_k\,x^{(k+1)}(t) + \frac{c_n}{\tau} \left(x^{(n)}(t+\tau)-x^{(n)}(t)\right)
\end{eqnarray*}
This is equivalent to:
\[
	x^{(n)}(t+\tau)
	\approx x^{(n)}(t) - \frac{\tau}{c_n}\,\sum_{k=0}^{n-1} c_k\,x^{(k+1)}(t)
\]
From this, we can read off the desired transition matrix $W$ of size $(n+1) \times (n+1)$ from \cref{mat}.
Together with a start vector $s$ satisfying \cref{diff}, we can thus solve
differential equations approximately by LRNNs.

\section{Proof of \cref{chad}}\label{proof:chad}

We first prove the case where the Jordan matrix $J$ only contains ordinary
Jordan blocks as in \cref{jordan}, i.e., possibly with complex eigenvalues
on the diagonal. Since $J$ is a direct sum of Jordan blocks, it suffices to
consider the case where $J$ is a single Jordan block because, as the Jordan
matrix $J$, the matrices $A$ and also $B$ (see below) can be obtained as direct
sums, too.

In the following, we use the column vectors $y = \big[ y_1 \cdots y_N
\big]^\top$ with all non-zero entries, $x = \big[ x_1 \cdots x_N \big]^\top$
with $x = V^{-1} \cdot s$ (cf. \cref{eigen}), and $b = \big[ b_1 \cdots b_N
\big]^\top$. From $b$, we construct the following upper triangular Toeplitz
matrix
\[ B = \left[ \begin{array}{*{4}{c}}
  b_N & \cdots & b_2 & b_1\\
  0 & \ddots & & b_2\\
  \vdots & \ddots & \ddots & \vdots\\
  0 & \cdots & 0 & b_N
\end{array} \right] \]
which commutes with the Jordan block $J$ \cite[Sect.~3.2.4]{HJ13}, i.e., it
holds that (a)~$J \cdot B = B \cdot J$. We define $B$ and hence $b$ by the
equation (b)~$x = B \cdot y$ which is equivalent to:

\[ \left[ \begin{array}{*{4}{c}}
  y_N & \cdots & y_2 & y_1\\
  0 & \ddots & & y_2\\
  \vdots & \ddots & \ddots & \vdots\\
  0 & \cdots & 0 & y_N
\end{array} \right] \cdot b = \left[ \begin{array}{c}
  x_1\\[7pt] x_2\\ \vdots\\ x_N
\end{array} \right] \]
Since the main diagonal of the left matrix contains no $0$s because $y_N \neq 0$
by precondition, there always exists a solution for $b$ \cite[Sect.~0.9.3]{HJ13}.
Then $A = V \cdot B$ does the job:
\[ f(t) = W^t \cdot s
	\overset{\text{\cref{jordan}}}{=} V \cdot J^t \cdot V^{-1} \cdot s
	= V \cdot J^t \cdot x
	\overset{\text{(b)}}{=} V \cdot J^t \cdot B \cdot y
	\overset{\text{(a)}}{=} V \cdot B \cdot J^t \cdot y
	= A \cdot J^t \cdot y
\]

The generalization to the real Jordan decomposition is straightforward by
applying the fact that for complex conjugate eigenvalue pairs $\lambda$ and
$\overline{\lambda}$ the matrix $M$ in a real Jordan block (cf. \cref{real}) is
similar to the diagonal matrix $D = \left[ \begin{array}{cc}
	\lambda & 0\\
	0 & \overline{\lambda}
\end{array} \right]$ via $U = \left[ \begin{array}{cc}
	-\mathfrak{i} & -\mathfrak{i}\\
	1 & -1
\end{array} \right]$ \citep[Sect.~3.4.1]{HJ13}, i.e., $M = U \cdot D \cdot
U^{-1}$. The above-mentioned commutation property~(a) analogously holds for real
Jordan blocks. This completes the proof.

\section{Proof of \cref{exponential}}\label{proof:exponential}

Let $f_k(t)$ denote the value of the $k$-th dimension of $f(t)$, $\lambda$ be
the eigenvalue of $W$ with maximal absolute value and $m$ be the maximal (geometric)
multiplicity of the eigenvalues of the transition matrix $W$. Then, from
\cref{jordan}, we can easily deduce \[ |f_k(t)| = O(t^m\,|\lambda|^t) \]
as asymptotic behavior for large $t$.

\section{Proof of \cref{approx}}\label{proof:approx}

First, we take the series of function values $f(t_0),\dots,f(t_n)$ and identify
them with the time series $S(0),\dots,S(n)$. After applying the LRNN learning
procedure, the LRNN runs through all given values, because by construction the
upper part $\big[ S(0) \cdots S(n-1) \big]$ of the matrix $X$ (cf. \cref{Xin})
and the matrix $Y^\mathrm{out}$ (cf. \cref{Yout}) consist of the series of
function values of $f$, provided that the linear
matrix equation $Y^\mathrm{out} = W^\mathrm{out} \cdot X$ (\cref{linear}) has at
least one solution.

Since \cref{linear} is equivalent to simultaneously solving the equations $y_k =
w_k \cdot X$, for $1 \le k \le N$ where $y_1,\dots,y_N$ and $w_1,\dots,w_N$
denote the row vectors of the matrices $Y^\mathrm{out}$ and $W^\mathrm{out}$,
respectively, the latter is the case if the rank of the coefficient matrix $X$
is equal to the rank of the augmented matrix $M_k = \big[ X ~~ y_k \big]^\top$
for every $k$. This leads to the equation $\mathrm{rank}(X) =
\min(N^\mathrm{in\,out}+N^\mathrm{res},n) = \mathrm{rank}(M_k) =
\min(N^\mathrm{in\,out}+N^\mathrm{res}+1,n)$. From this, it follows that, as
desired, $N^\mathrm{res} \ge n-N^\mathrm{in\,out}$ reservoir neurons have to be
employed to guarantee at least one solution for the $w_k$, provided that the
rank of the matrix $X$ is maximal. For the latter, we consider two cases:
\begin{itemize}
  \item If the rank of the upper part of the matrix $X$ (see above) is not
	maximal, then this does not cause any problems. We only have to replace
	$N^\mathrm{in\,out}$ by the actual rank of the upper part of the matrix
	$X$ in \cref{ineq}.
  \item The rank of the lower part $\big[ R(0) \cdots R(n-1) \big]$ of the
	matrix $X$ almost always has maximal rank, because we employ a random
	reservoir (cf. \cref{thedef}). Thus, a suitable reservoir does the job,
	which completes the proof.
\end{itemize}


\end{document}
